\documentclass[platex]{jlreq}
\usepackage{url}

\title{AIでソフトウェアを開発できても、それが研究の目的に合い、結果を信頼できるものか確かめることが本質的なボトルネックになるという話\par
       \large SPReADコミュニティ発表の参考文献リスト}
\author{発表者\quad 山崎 進}
\date{}

\begin{document}
\maketitle

発表を視聴した方が，ソフトウェアの要求（requirements）と
妥当性確認（validation）を学ぶための文献を，日英両方で読める2点に絞った。
ここでは，作る前に「誰が何のために使い，何ができる必要があるか」を
明らかにする段階と，できたものが「実際の利用目的を満たすか」を
確かめる段階に注目する。

\section*{要求を明らかにする}

\paragraph{『ソフトウェア要求 第3版』\cite{wiegers2014ja,wiegers2013}}
第1章で要求の基本を押さえ，第7章で利用者から必要なことを聞き出す方法，
第11章で確認できる形に書く方法を学べる。
第17章は，書いた要求が利用者の目的に合っているかを
開発前に見直し，完成後の合否を判断する受け入れ基準につなげる章である。
たとえば「測定値を集計したい」という要望なら，
その集計で何を判断するのかを聞き，欠測値をどう扱うか，
どの件数と平均を出すかを具体的に決めておく。

\section*{できたものが利用目的を満たすか確かめる}

\paragraph{『Foundation Level シラバス』\cite{jstqb2023,istqb2024}}
日英とも無料で読める。第1章の1.1節で，
書かれた要求を満たすかを調べる検証（verification）と，
利用者の必要を満たすかを調べる妥当性確認の違いを押さえ，
第2章の2.2.1節にある受け入れテストの説明へ進むとよい。
受け入れテストでは，利用者が実際の使い方を想定して，
その目的に使えるかを確かめる。
たとえば，想定する条件を含む測定データで集計を行い，
件数と平均を独立に求めた値と照合したうえで，
その集計が研究上の判断に必要な情報を提供するかを確認する。

\bibliographystyle{jname}
\bibliography{ref}

\end{document}
